\documentclass[../../uem_phy_std_q.tex]{subfiles}

\begin{document}
    \section{静電場}
    \subsection{点電荷が作る電場}
        \subfile{01_point-charge/uem_phy_std_em_ef_01_q.tex}
        \vspace{20pt}
        \vfill
    \subsection{平行一様電場・電位の基準}
        \subfile{02_uniform/uem_phy_std_em_ef_02_q.tex}
        \vspace{20pt}
        \vfill
    \subsection{拡がった電荷分布の作る電場}
        \subfile{03_continuous-charge/uem_phy_std_em_ef_03_q.tex}
        \vspace{20pt}
        \vfill
    \subsection{電位から電場を逆算する①}
        \subfile{04_from-potential-1/uem_phy_std_em_ef_04_q.tex}
        \vspace{20pt}
        \vfill
    \subsection{電位から電場を逆算する②}
        \subfile{05_from-potential-2/uem_phy_std_em_ef_05_q.tex}
        \vspace{20pt}
        \vfill
    \subsection{※点電荷が作る電場}
        \subfile{06_add-point-charge/uem_phy_std_em_ef_06_q.tex}
        \vspace{20pt}
        \vfill
    \subsection{※平行一様電場}
        \subfile{07_add-uniform/uem_phy_std_em_ef_07_q.tex}
        \vspace{20pt}
        \vfill
    \subsection{※電気力線}
        \subfile{08_add-field-lines/uem_phy_std_em_ef_08_q.tex}
        \vspace{20pt}
        \vfill
    \subsection{※等電位線}
        \subfile{09_add-equipotential-lines/uem_phy_std_em_ef_09_q.tex}
        \vspace{20pt}
        \vfill
\end{document}
